\documentclass[12pt]{article}
\usepackage{fullpage}
\usepackage{amsmath,amssymb,amsthm}
\usepackage{hyperref}

\newtheorem{theorem}{Theorem}
\newtheorem{lemma}{Lemma}
\newtheorem{proposition}{Proposition}
\newtheorem{corollary}{Corollary}
\theoremstyle{definition}
\newtheorem{definition}{Definition}
\newtheorem{remark}{Remark}

\DeclareMathOperator{\sgn}{sgn}

\begin{document}

\begin{center}
\textbf{\large A rationally generating family of $4k$--manifolds, by direct linear
algebra}\\[2pt]
\textit{(self-contained resolution of the ``generation gap'', issue rm-01541-2)}
\end{center}

\bigskip

\noindent
\textbf{Scope.}
This section supplies, by elementary means, the linear-algebra step that the standard
proof of the Hirzebruch signature theorem usually imports from cobordism theory
(Thom's computation $\Omega^{SO}_\ast\otimes\mathbb{Q}=
\mathbb{Q}[[\mathbb{CP}^2],[\mathbb{CP}^4],\dots]$).  We never use the cobordism ring,
the Pontryagin--Thom construction, or any spectral/index-theoretic input.  Everything
below is a finite computation with Pontryagin classes of explicit manifolds and a
determinant estimate.

\medskip

\noindent
\textbf{What is proved here.}
For each $k\ge 1$ we exhibit a finite, explicit family of closed oriented
$4k$--manifolds whose Pontryagin-number vectors are linearly independent and span the
full space of Pontryagin numbers in degree $4k$ (Theorem~\ref{thm:span}); we compute in
closed form the value on each family member of both functionals occurring in the
signature theorem, the signature $\tau$ and the $L$-number
$L[M]=\langle L(p_\bullet(M)),[M]\rangle$ (Proposition~\ref{prop:values}); and we deduce
that these two functionals coincide on degree-$4k$ Pontryagin numbers
(Corollary~\ref{cor:agree}).  The only external inputs are: (i) the multiplicativity of
both $\tau$ and $L[\cdot]$ under products, and (ii) the fact (established elsewhere in
the proof) that $\tau$ is expressible through Pontryagin numbers in each degree.  Both
(i) and (ii) are recalled precisely below; the genuinely combinatorial content -- the
spanning -- is proved here from scratch.

\bigskip

\section{The space of Pontryagin numbers and the two functionals}

Throughout, $k\ge 1$ is fixed and $P(k)$ denotes the set of partitions of $k$; we write
$p(k)=\#P(k)$.  A partition is written $\lambda=(\lambda_1\ge\lambda_2\ge\cdots\ge
\lambda_r\ge 1)$ with $|\lambda|=\sum_i\lambda_i=k$ and length $\ell(\lambda)=r$.

\begin{definition}[Pontryagin numbers]\label{def:pn}
Let $M$ be a closed oriented smooth $4k$--manifold with Pontryagin classes
$p_i=p_i(TM)\in H^{4i}(M;\mathbb{Q})$.  For $\lambda\in P(k)$ put
\[
p_\lambda[M]\;:=\;\big\langle\, p_{\lambda_1}\smile p_{\lambda_2}\smile\cdots\smile
p_{\lambda_r}\,,\,[M]\,\big\rangle\;\in\;\mathbb{Q}.
\]
The \emph{Pontryagin vector} of $M$ is
$\mathbf{p}[M]=\big(p_\lambda[M]\big)_{\lambda\in P(k)}\in\mathbb{Q}^{P(k)}$.
\end{definition}

The $p(k)$ monomials $p_\lambda$ form a $\mathbb{Q}$-basis of the degree-$4k$ part of
the free graded-commutative polynomial algebra $\mathbb{Q}[p_1,p_2,\dots]$ on generators
$p_i$ of degree $4i$.  Consequently the degree-$4k$ component of any multiplicative
sequence (in particular Hirzebruch's $L$-polynomial) is a $\mathbb{Q}$-linear
combination of the $p_\lambda$; the $L$-number is therefore a fixed linear functional of
the Pontryagin vector:
\[
L[M]\;=\;\big\langle L(p_\bullet(M)),[M]\big\rangle\;=\;
\sum_{\lambda\in P(k)} c_\lambda\, p_\lambda[M],
\qquad c_\lambda\in\mathbb{Q}\ \text{fixed},
\tag{1}\label{eq:Llinear}
\]
where the constants $c_\lambda$ are the coefficients of the degree-$4k$ part of the
$L$-series (the precise values of the $c_\lambda$ are irrelevant for the present step).

The signature theorem is the assertion that the integer-valued functional $M\mapsto
\tau(M)$ equals the functional \eqref{eq:Llinear}.  In the elementary proof the
following two facts are established by independent (non-cobordism) arguments, and we take
them as the interface to the rest of the proof:

\begin{itemize}
\item[\textbf{(I)}] \emph{Multiplicativity.} For closed oriented manifolds $M,N$ of
dimensions divisible by $4$,
\[
\tau(M\times N)=\tau(M)\,\tau(N),\qquad L[M\times N]=L[M]\cdot L[N].
\tag{2}\label{eq:mult}
\]
(For $L[\cdot]$ this is the multiplicativity of multiplicative sequences together with
the Whitney formula $p(T(M\times N))=p(TM)\times p(TN)$ and
$[M\times N]=[M]\times[N]$; for $\tau$ it is the elementary fact that the intersection
form of a product is the tensor product of the intersection forms, whose signature is
the product of the signatures.)
\item[\textbf{(II)}] \emph{$\tau$ is a Pontryagin functional.} In each degree $4k$ there
exist (a priori unknown) constants $a_\lambda\in\mathbb{Q}$ with
$\tau(M)=\sum_{\lambda\in P(k)}a_\lambda\, p_\lambda[M]$ for every closed oriented
$4k$--manifold $M$.
\end{itemize}

\noindent
Fact \eqref{eq:mult} is elementary and recalled above.  Fact (II) is the
non-combinatorial heart of the theorem and is supplied elsewhere in the proof; it does
\emph{not} rely on the present section.  The role of the present section is to show that
(I), (II), and the value computations below force $a_\lambda=c_\lambda$ for all
$\lambda$, i.e.\ $\tau=L[\cdot]$.

\bigskip

\section{The generating family and its Pontryagin numbers}

\begin{definition}[The family]\label{def:family}
For $\mu=(\mu_1\ge\cdots\ge\mu_r)\in P(k)$ set
\[
M_\mu\;:=\;\mathbb{CP}^{2\mu_1}\times\mathbb{CP}^{2\mu_2}\times\cdots\times
\mathbb{CP}^{2\mu_r},
\]
a closed oriented manifold of real dimension $\sum_s 4\mu_s=4k$.  This is a family of
exactly $p(k)$ manifolds, indexed by $P(k)$.
\end{definition}

We now compute all Pontryagin numbers of $M_\mu$ in closed form.  Recall the standard
fact, proved purely from the Euler sequence and the splitting principle: with
$h\in H^2(\mathbb{CP}^n;\mathbb{Z})$ the hyperplane class,
\[
p(T\mathbb{CP}^n)\;=\;(1+h^2)^{\,n+1},
\qquad\text{so}\qquad
p_i(\mathbb{CP}^n)=\binom{n+1}{i}h^{2i},
\tag{3}\label{eq:cppont}
\]
and $\langle h^n,[\mathbb{CP}^n]\rangle=1$.  (Equation \eqref{eq:cppont} is elementary:
$T\mathbb{CP}^n\oplus\underline{\mathbb{C}}\cong(n+1)\,\overline{\mathcal{O}(1)}$, whence
$c(T\mathbb{CP}^n)=(1+h)^{n+1}$ and, in the splitting calculus, the total Pontryagin
class is obtained by replacing each Chern root $x_j$ with $x_j^2$, giving
$(1+h^2)^{n+1}$ since all Chern roots equal $h$.)

Introduce, for each factor $s$ of $M_\mu$, the degree-$4$ class
\[
t_s\;:=\;h_s^2\in H^4(M_\mu;\mathbb{Q}),
\qquad s=1,\dots,r,
\]
where $h_s$ is the hyperplane class pulled back from the $s$-th factor
$\mathbb{CP}^{2\mu_s}$.  By \eqref{eq:cppont} and the Whitney formula, the total
Pontryagin class of $M_\mu$ is
\[
p(TM_\mu)\;=\;\prod_{s=1}^{r}(1+t_s)^{\,2\mu_s+1},
\tag{4}\label{eq:prodp}
\]
and the fundamental cohomology class is dual to $\prod_s h_s^{2\mu_s}=\prod_s t_s^{\mu_s}$,
so for any polynomial $Q$ in the $t_s$,
\[
\langle Q,[M_\mu]\rangle\;=\;\big[\textstyle\prod_s t_s^{\mu_s}\big]\,Q,
\tag{5}\label{eq:pair}
\]
the coefficient of the monomial $\prod_s t_s^{\mu_s}$ in $Q$.

It is convenient to work with the \emph{power-sum Pontryagin numbers}.  Let
\[
\mathfrak{p}_n\;:=\;\sum_{j} y_j^{\,n}
\]
denote the $n$-th power sum of the formal Pontryagin roots $y_j$ (so that, as usual,
$1+p_1+p_2+\cdots=\prod_j(1+y_j)$ and the $\mathfrak{p}_n$ are the Newton polynomials in
the $p_i$ with \emph{rational} coefficients).  For a partition
$\lambda=(\lambda_1,\dots,\lambda_r)\in P(k)$ define
\[
\mathfrak{p}_\lambda\;:=\;\mathfrak{p}_{\lambda_1}\mathfrak{p}_{\lambda_2}\cdots
\mathfrak{p}_{\lambda_r},
\qquad
\mathfrak{p}_\lambda[M]\;:=\;\langle\mathfrak{p}_\lambda,[M]\rangle.
\]
Because the Newton relations express the $\mathfrak{p}_n$ in terms of the $p_i$ and vice
versa by an \emph{invertible} (unitriangular up to nonzero scalars) change of basis over
$\mathbb{Q}$, the families $\{\mathfrak{p}_\lambda\}_{\lambda\in P(k)}$ and
$\{p_\lambda\}_{\lambda\in P(k)}$ span the same degree-$4k$ subspace and
\[
\operatorname{span}_\mathbb{Q}\{\mathbf{p}[M_\mu]:\mu\in P(k)\}
=\mathbb{Q}^{P(k)}
\iff
\det\big(\mathfrak{p}_\lambda[M_\mu]\big)_{\lambda,\mu\in P(k)}\neq 0.
\tag{6}\label{eq:reduce}
\]

For $M_\mu$, the formal Pontryagin roots are, by \eqref{eq:prodp}, the classes $t_s$ each
occurring with multiplicity $2\mu_s+1$.  Hence
\[
\mathfrak{p}_n(M_\mu)=\sum_{s=1}^{r}(2\mu_s+1)\,t_s^{\,n},
\]
and therefore, using \eqref{eq:pair},
\[
\mathfrak{p}_\lambda[M_\mu]
=\Big[\textstyle\prod_{s} t_s^{\mu_s}\Big]\
\prod_{i=1}^{\ell(\lambda)}\Big(\sum_{s=1}^{\ell(\mu)}(2\mu_s+1)\,t_s^{\lambda_i}\Big).
\tag{7}\label{eq:closed}
\]
Formula \eqref{eq:closed} is an explicit, finite, integer-valued closed form for every
entry of the matrix in \eqref{eq:reduce}.

\bigskip

\section{Triangularity and nonsingularity}

We now prove that the matrix
$A:=\big(\mathfrak{p}_\lambda[M_\mu]\big)_{\lambda,\mu\in P(k)}$ is nonsingular, by a
direct combinatorial argument (no cobordism, no symmetric-function black boxes beyond the
elementary Newton change of basis already invoked).

Recall the \emph{dominance order} on $P(k)$: for $\alpha,\beta\in P(k)$ written in
weakly decreasing order and padded with zeros to equal length,
\[
\alpha\preceq\beta\iff
\sum_{i\le j}\alpha_i\le\sum_{i\le j}\beta_i\quad\text{for all }j.
\]

\begin{lemma}[Combinatorial expansion of the entries]\label{lem:expand}
For $\lambda,\mu\in P(k)$,
\[
\mathfrak{p}_\lambda[M_\mu]
=\sum_{f}\ \prod_{i=1}^{\ell(\lambda)}\big(2\mu_{f(i)}+1\big),
\tag{8}\label{eq:expand}
\]
where the sum runs over all functions $f:\{1,\dots,\ell(\lambda)\}\to\{1,\dots,\ell(\mu)\}$
such that
\[
\sum_{i:\,f(i)=s}\lambda_i=\mu_s\quad\text{for every }s=1,\dots,\ell(\mu).
\tag{9}\label{eq:balance}
\]
In particular every entry $\mathfrak{p}_\lambda[M_\mu]$ is a nonnegative integer, and it
is nonzero if and only if at least one such function $f$ exists.
\end{lemma}

\begin{proof}
Expand the product in \eqref{eq:closed}.  Each factor indexed by $i$ contributes one term
$(2\mu_{s}+1)\,t_{s}^{\lambda_i}$, i.e.\ a choice $f(i)=s\in\{1,\dots,\ell(\mu)\}$ of which
summand is taken.  A choice of $f$ produces the monomial
$\prod_{i}(2\mu_{f(i)}+1)\,t_{f(i)}^{\lambda_i}
=\big(\prod_i(2\mu_{f(i)}+1)\big)\,\prod_{s}t_s^{\sum_{i:f(i)=s}\lambda_i}$.
By \eqref{eq:pair} the coefficient of $\prod_s t_s^{\mu_s}$ is the sum of
$\prod_i(2\mu_{f(i)}+1)$ over precisely those $f$ for which the exponent of each $t_s$
equals $\mu_s$, which is condition \eqref{eq:balance}.  All weights
$2\mu_s+1\ge 3>0$, so the sum is a nonnegative integer, zero exactly when no $f$ satisfies
\eqref{eq:balance}.
\end{proof}

\begin{lemma}[Coarsening dominates]\label{lem:merge}
Let $\lambda\in P(k)$ and let $\{B_1,\dots,B_m\}$ be a partition of the index set
$\{1,\dots,\ell(\lambda)\}$ into nonempty blocks.  Define $\mu_s:=\sum_{i\in B_s}\lambda_i$
for $s=1,\dots,m$; then $\mu=(\mu_1,\dots,\mu_m)$ is a partition of $k$ (after weakly
decreasing rearrangement) and $\lambda\preceq\mu$ in the dominance order.
\end{lemma}

\begin{proof}
First, a single \emph{pairwise merge}.  Let $\nu$ be a partition and let $a\ge b\ge 1$ be
two of its parts; let $\nu'$ be obtained by deleting one occurrence each of $a$ and $b$ and
inserting the single part $a+b$ (so $|\nu'|=|\nu|$).  We claim $\nu\preceq\nu'$.  Write both
partitions in weakly decreasing order.  Passing from $\nu$ to $\nu'$ replaces two parts
$a,b$ by the larger part $a+b$ and one fewer total part.  Consider the partial sums
$S_j(\cdot):=\sum_{i\le j}(\cdot)^{\downarrow}_i$.  For any threshold value, the multiset of
parts of $\nu'$ that are $\ge$ that threshold is obtained from that of $\nu$ by possibly
removing $a$ and $b$ and possibly adding $a+b\ge a\ge b$.  Concretely: the part $a+b$ in
$\nu'$ is at least as large as either of $a,b$, so for every $j$ the $j$ largest parts of
$\nu'$ have sum at least the $j$ largest parts of $\nu$.  Indeed, fix $j$.  Let $T$ be the
multiset of the $j$ largest parts of $\nu$ (sum $S_j(\nu)$).  Build a sub-multiset $T'$ of
the parts of $\nu'$ with $|T'|\le j$ and $\sum T'\ge \sum T$ as follows: keep every element
of $T$ that is still a part of $\nu'$ (i.e.\ all elements of $T$ except possibly the deleted
copies of $a$ and $b$); if $T$ contained at least one of $a,b$, replace those (one or two)
deleted elements by the single part $a+b$, which is $\ge$ their sum if both were present and
$\ge$ the single one otherwise.  This $T'$ has $|T'|\le j$ and $\sum T'\ge\sum T=S_j(\nu)$.
Since the $j$ largest parts of $\nu'$ maximize the sum over all sub-multisets of size
$\le j$, we get $S_j(\nu')\ge\sum T'\ge S_j(\nu)$.  As $j$ was arbitrary and total sums
agree, $\nu\preceq\nu'$.

Now the general statement.  The passage from $\lambda$ to $\mu$ replaces each block
$B_s=\{i_1,\dots,i_{|B_s|}\}$ by the single part $\sum_{i\in B_s}\lambda_i$, which is the
result of $|B_s|-1$ successive pairwise merges within that block; performing all such merges
across all blocks expresses $\mu$ as the end of a finite chain
$\lambda=\nu^{(0)}\preceq\nu^{(1)}\preceq\cdots\preceq\nu^{(N)}=\mu$, each step a pairwise
merge.  By the previous paragraph each step is a dominance increase, and $\preceq$ is
transitive, so $\lambda\preceq\mu$.
\end{proof}

\begin{lemma}[Triangularity]\label{lem:tri}
If $\mathfrak{p}_\lambda[M_\mu]\neq 0$ then $\lambda\preceq\mu$.  Moreover
$\mathfrak{p}_\mu[M_\mu]>0$ for every $\mu\in P(k)$.
\end{lemma}

\begin{proof}
Suppose $\mathfrak{p}_\lambda[M_\mu]\neq 0$.  By Lemma~\ref{lem:expand} there is a
function $f:\{1,\dots,\ell(\lambda)\}\to\{1,\dots,\ell(\mu)\}$ satisfying
\eqref{eq:balance}.  The fibers $B_s:=f^{-1}(s)$ partition $\{1,\dots,\ell(\lambda)\}$; by
\eqref{eq:balance} each is nonempty (if some $B_s$ were empty then $\mu_s=0$, contradicting
$\mu_s\ge 1$), and $\mu_s=\sum_{i\in B_s}\lambda_i$.  Thus $\mu$ is the coarsening of
$\lambda$ associated to the blocks $\{B_s\}$, so by Lemma~\ref{lem:merge},
$\lambda\preceq\mu$.

For the diagonal, take $\lambda=\mu$ and let $f=\mathrm{id}$ (after fixing any ordering of
the parts of $\mu$); then $\sum_{i:f(i)=s}\mu_i=\mu_s$, so $f$ satisfies \eqref{eq:balance}.
By \eqref{eq:expand} the entry $\mathfrak{p}_\mu[M_\mu]$ is a sum of strictly positive terms
(one for each such $f$), hence $\mathfrak{p}_\mu[M_\mu]>0$.
\end{proof}

\begin{theorem}[Spanning]\label{thm:span}
For every $k\ge 1$ the $p(k)$ manifolds $\{M_\mu:\mu\in P(k)\}$ have linearly independent
Pontryagin vectors, and these vectors span $\mathbb{Q}^{P(k)}$.  Equivalently, the only
$\mathbb{Q}$-linear functional of degree-$4k$ Pontryagin numbers that vanishes on every
$M_\mu$ is the zero functional.
\end{theorem}

\begin{proof}
Order $P(k)$ by any linear extension of the dominance order $\preceq$ (such an extension
exists since $\preceq$ is a partial order on the finite set $P(k)$).  By
Lemma~\ref{lem:tri}, in this ordering the matrix
$A=\big(\mathfrak{p}_\lambda[M_\mu]\big)_{\lambda,\mu}$ is triangular: its $(\lambda,\mu)$
entry vanishes whenever $\lambda\not\preceq\mu$, i.e.\ whenever $\lambda$ comes strictly
after $\mu$ in the chosen ordering.  Its diagonal entries
$\mathfrak{p}_\mu[M_\mu]>0$ are nonzero.  Hence
\[
\det A=\prod_{\mu\in P(k)}\mathfrak{p}_\mu[M_\mu]\;>\;0,
\]
so $A$ is nonsingular.  By the equivalence \eqref{eq:reduce}, the Pontryagin vectors
$\{\mathbf{p}[M_\mu]\}$ are linearly independent and span $\mathbb{Q}^{P(k)}$.  The final
sentence is the dual statement: a linear functional vanishing on a spanning set is zero.
\end{proof}

\begin{remark}
The determinant is in fact a product of positive integers; for small $k$ one finds
$\det A=3,\,90,\,17010,\,1653372000,\dots$ for $k=1,2,3,4$.  The argument uses only the
explicit Pontryagin classes \eqref{eq:cppont}, the Whitney formula, the elementary Newton
change of basis between $p_\lambda$ and $\mathfrak{p}_\lambda$, and the combinatorial
triangularity of Lemma~\ref{lem:tri}.  At no point is $\Omega^{SO}_\ast\otimes\mathbb{Q}$,
the Pontryagin--Thom construction, or any generation statement from cobordism theory
invoked; the spanning is a self-contained determinant computation.
\end{remark}

\bigskip

\section{Closed-form values of the two functionals}

\begin{proposition}[Values on the family]\label{prop:values}
For every $k\ge 1$ and every $\mu\in P(k)$,
\[
\tau(M_\mu)=1
\qquad\text{and}\qquad
L[M_\mu]=1.
\]
\end{proposition}

\begin{proof}
\emph{Signature.}  For even $m$, the manifold $\mathbb{CP}^{m}$ has middle cohomology
$H^{m}(\mathbb{CP}^{m};\mathbb{R})=\mathbb{R}\cdot h^{m/2}$ one-dimensional, with
intersection form $(h^{m/2},h^{m/2})\mapsto\langle h^{m},[\mathbb{CP}^{m}]\rangle=1$.
This form is positive definite, so $\tau(\mathbb{CP}^{m})=1$ for even $m$; in particular
$\tau(\mathbb{CP}^{2\mu_s})=1$ for each $s$.  By the multiplicativity
\eqref{eq:mult}, $\tau(M_\mu)=\prod_s\tau(\mathbb{CP}^{2\mu_s})=1$.

\emph{$L$-number.}  By multiplicativity \eqref{eq:mult} of the $L$-genus it suffices to
show $L[\mathbb{CP}^{2m}]=1$ for all $m\ge 1$.  By \eqref{eq:cppont} the formal
Pontryagin roots of $\mathbb{CP}^{2m}$ are $h^2$ with multiplicity $2m+1$; hence the
$L$-polynomial evaluated on $\mathbb{CP}^{2m}$ is the degree-$4m$ part of
\[
L(p_\bullet)\Big|_{\mathbb{CP}^{2m}}
=\Big(\frac{\sqrt{h^2}}{\tanh\sqrt{h^2}}\Big)^{2m+1}
=\Big(\frac{h}{\tanh h}\Big)^{2m+1},
\]
where $h/\tanh h=\sum_{j\ge 0}c_j h^{2j}$ is the even power series defining the $L$-genus.
Pairing with $[\mathbb{CP}^{2m}]$ extracts the coefficient of $h^{2m}$:
\[
L[\mathbb{CP}^{2m}]=\big[h^{2m}\big]\Big(\frac{h}{\tanh h}\Big)^{2m+1}
=\frac{1}{2\pi i}\oint \Big(\frac{h}{\tanh h}\Big)^{2m+1}\frac{dh}{h^{2m+1}}
=\frac{1}{2\pi i}\oint\frac{dh}{(\tanh h)^{2m+1}},
\]
the contour being a small circle about $h=0$.  Substitute $w=\tanh h$, so
$dw=(1-w^2)\,dh=\operatorname{sech}^2 h\,dh$ and $h=0\mapsto w=0$ with $dw/dh=1$ at the
origin; the substitution is a local biholomorphism near $0$, and
\[
\frac{1}{2\pi i}\oint\frac{dh}{(\tanh h)^{2m+1}}
=\frac{1}{2\pi i}\oint\frac{1}{w^{2m+1}}\,\frac{dw}{1-w^2}
=\big[w^{2m}\big]\frac{1}{1-w^2}
=\big[w^{2m}\big]\sum_{j\ge0}w^{2j}=1.
\]
Hence $L[\mathbb{CP}^{2m}]=1$, and by multiplicativity
$L[M_\mu]=\prod_s L[\mathbb{CP}^{2\mu_s}]=1$.
\end{proof}

\begin{corollary}[Agreement on Pontryagin numbers]\label{cor:agree}
For every $k\ge 1$, the two $\mathbb{Q}$-linear functionals of degree-$4k$ Pontryagin
numbers,
\[
M\longmapsto\tau(M)
\qquad\text{and}\qquad
M\longmapsto L[M]=\sum_\lambda c_\lambda\,p_\lambda[M],
\]
coincide; that is, $a_\lambda=c_\lambda$ for all $\lambda\in P(k)$, where $a_\lambda$ are
the constants of fact \textnormal{(II)}.  Consequently
\[
\tau(M)=L[M]=\big\langle L(p_\bullet(M)),[M]\big\rangle
\]
for every closed oriented smooth $4k$--manifold $M$.
\end{corollary}

\begin{proof}
By fact (II), $\tau(M)=\sum_\lambda a_\lambda p_\lambda[M]$ for fixed $a_\lambda$, while by
\eqref{eq:Llinear} $L[M]=\sum_\lambda c_\lambda p_\lambda[M]$.  Their difference
\[
\Delta(M):=\tau(M)-L[M]=\sum_{\lambda\in P(k)}(a_\lambda-c_\lambda)\,p_\lambda[M]
\]
is a $\mathbb{Q}$-linear functional of the degree-$4k$ Pontryagin numbers.  By
Proposition~\ref{prop:values}, $\Delta(M_\mu)=\tau(M_\mu)-L[M_\mu]=1-1=0$ for every
$\mu\in P(k)$.  Thus $\Delta$ vanishes on the spanning family of
Theorem~\ref{thm:span}, hence $\Delta\equiv 0$ as a functional, i.e.\
$a_\lambda=c_\lambda$ for all $\lambda$.  Therefore $\tau(M)=L[M]$ for every closed
oriented $4k$--manifold $M$.
\end{proof}

\bigskip

\begin{remark}[On the interface facts]
The corollary closes the signature theorem \emph{given} facts (I) and (II).  Fact (I) is
elementary and proved above.  Fact (II) -- that $\tau$ is, in each degree, a fixed linear
combination of Pontryagin numbers -- is the analytic/topological core of the theorem and
is established in the companion sections of this proof by elementary (non-cobordism)
means; it is logically independent of the present section.  The content delivered here,
and the precise object the cobordism-free argument was missing, is the
\emph{spanning family with closed-form values}: an explicit set of $p(k)$ manifolds whose
Pontryagin vectors are provably independent by a direct triangular-determinant
computation, on which both functionals are computed to equal $1$.  This is exactly what
the standard proof imports from $\Omega^{SO}_\ast\otimes\mathbb{Q}$, now supplied
elementarily.
\end{remark}

\end{document}
